\documentclass[UTF8]{ctexart}
\usepackage{amsmath}
\usepackage{tocloft}
\usepackage[bigfiles]{pdfbase}
\usepackage{amsthm,amssymb}
\usepackage{booktabs}
\usepackage{graphicx}
\usepackage[hidelinks]{hyperref}
\usepackage{geometry}
\usepackage{float}
\usepackage{multirow}
\usepackage{indentfirst}
\usepackage{diagbox}
\usepackage{listings}
\usepackage{xcolor}
\usepackage{markdown}
\definecolor{codegreen}{rgb}{0,0.6,0}
\definecolor{codegray}{rgb}{0.5,0.5,0.5}
\definecolor{codepurple}{rgb}{0.58,0,0.82}
\definecolor{backcolour}{rgb}{0.95,0.95,0.92}

\lstdefinestyle{mystyle}{
    backgroundcolor=\color{backcolour},   
    commentstyle=\color{codegreen},
    keywordstyle=\color{magenta},
    numberstyle=\tiny\color{codegray},
    stringstyle=\color{codepurple},
    basicstyle=\ttfamily\footnotesize,
    breakatwhitespace=false,         
    breaklines=true,                 
    captionpos=b,                    
    keepspaces=true,                 
    numbers=left,                    
    numbersep=5pt,                  
    showspaces=false,                
    showstringspaces=false,
    showtabs=false,                  
    tabsize=2
}
\lstset{style=mystyle}
\newcommand\question[2]{\noindent\fbox{\parbox[t]{\linewidth}{\hspace{0pt}{\textbf{#1\quad}#2}}}}


\usepackage{graphicx}
\usepackage{setspace}
\renewcommand{\cftsecleader}{\cftdotfill{\cftdotsep}}
\geometry{a4paper,left=3cm,right=3cm,top=2cm,bottom=2cm} 


\CTEXsetup[format={\Large \bfseries}]{section}
\title{\textbf{人工智能赋能的数字集成电路设计 作业三}}
\author{陈天乐\ \ [student ID removed] }
\date{\today}
\begin{document}
\maketitle
\tableofcontents
\newpage

\section{功能一：代码补全 }
将代码选中，右键选择Generate Code的功能。

\begin{figure}[H]
    \centering
    \includegraphics[width=\textwidth]{代码补全.png}
    \caption{代码补全}
\end{figure}

生成的代码如下：
\begin{lstlisting}[frame=shadowbox,caption=Demo\_RTL\_Complete.v]
/*
This is a multiplier for floating point
- Support IEEE 754 Half-Precision Floating-Point Format (FP16):
	- Handle all special cases (zero, infinity, NaN) as specified by IEEE 754.
	- Ensure correct rounding to the nearest even value.
- Use valid and ready signals for handshake. When valid is high, the input operands A and B are ready. When ready is high, the result C is valid.
*/

module fp_multiplier (
    input clk,
    input rst_n,
    input [15:0] A,
    input [15:0] B,
    input Ready,
    output reg [15:0] C,
    output reg Valid
);
// Internal registers for pipeline stages
    reg [15:0] A_reg, B_reg;
    reg ready_reg;
    reg [1:0] state;
    
    // Floating point components extraction
    wire sign_A, sign_B;
    wire [4:0] exp_A, exp_B;
    wire [9:0] frac_A, frac_B;
    
    // Stage 1 signals
    reg stage1_valid;
    reg sign_prod;
    reg [5:0] exp_sum;
    reg [21:0] frac_prod;
    reg is_zero_A, is_zero_B;
    reg is_inf_A, is_inf_B;
    reg is_nan_A, is_nan_B;
    
    // Stage 2 signals
    reg stage2_valid;
    reg [4:0] exp_adj;
    reg [10:0] frac_rounded;
    reg is_zero_prod, is_inf_prod, is_nan_prod;
    
    // Constants
    localparam [4:0] EXP_BIAS = 5'd15;
    localparam [4:0] EXP_MAX = 5'd31;
    localparam [4:0] EXP_MIN = 5'd0;
    
    // Extract components
    assign sign_A = A_reg[15];
    assign exp_A = A_reg[14:10];
    assign frac_A = A_reg[9:0];
    
    assign sign_B = B_reg[15];
    assign exp_B = B_reg[14:10];
    assign frac_B = B_reg[9:0];
    
    // Special case detection
    always @(*) begin
        is_zero_A = (exp_A == 5'b0) && (frac_A == 10'b0);
        is_inf_A = (exp_A == 5'b11111) && (frac_A == 10'b0);
        is_nan_A = (exp_A == 5'b11111) && (frac_A != 10'b0);
        
        is_zero_B = (exp_B == 5'b0) && (frac_B == 10'b0);
        is_inf_B = (exp_B == 5'b11111) && (frac_B == 10'b0);
        is_nan_B = (exp_B == 5'b11111) && (frac_B != 10'b0);
    end
    
    // Pipeline stage 1: Multiplication and special case handling
    always @(posedge clk or negedge rst_n) begin
        if (!rst_n) begin
            stage1_valid <= 1'b0;
            sign_prod <= 1'b0;
            exp_sum <= 6'b0;
            frac_prod <= 22'b0;
        end else if (ready_reg) begin
            stage1_valid <= 1'b1;
            
            // Sign calculation
            sign_prod <= sign_A ^ sign_B;
            
            // Exponent calculation (subtract bias)
            exp_sum <= {1'b0, exp_A} + {1'b0, exp_B} - EXP_BIAS;
            
            // Fraction multiplication (with implicit 1)
            if (is_zero_A || is_zero_B) begin
                frac_prod <= 22'b0;
            end else begin
                frac_prod <= {1'b1, frac_A} * {1'b1, frac_B};
            end
        end else begin
            stage1_valid <= 1'b0;
        end
    end
    
    // Rounding logic (round to nearest even)
    wire [21:0] frac_normalized;
    wire [4:0] exp_normalized;
    wire round_bit, sticky_bit, guard_bit;
    wire round_up;
    
    // Normalization
    assign frac_normalized = (frac_prod[21]) ? (frac_prod >> 1) : frac_prod;
    assign exp_normalized = (frac_prod[21]) ? (exp_sum + 1) : exp_sum;
    
    // Rounding bits
    assign guard_bit = frac_normalized[10];
    assign round_bit = frac_normalized[9];
    assign sticky_bit = |frac_normalized[8:0];
    
    // Round to nearest even
    assign round_up = guard_bit & (round_bit | sticky_bit | frac_normalized[11]);
    
    // Pipeline stage 2: Normalization, rounding and final result assembly
    always @(posedge clk or negedge rst_n) begin
        if (!rst_n) begin
            stage2_valid <= 1'b0;
            C <= 16'b0;
            Valid <= 1'b0;
            exp_adj <= 5'b0;
            frac_rounded <= 11'b0;
            is_zero_prod <= 1'b0;
            is_inf_prod <= 1'b0;
            is_nan_prod <= 1'b0;
        end else if (stage1_valid) begin
            stage2_valid <= 1'b1;
            
            // Handle special cases
            is_nan_prod <= is_nan_A | is_nan_B | (is_inf_A & is_zero_B) | (is_zero_A & is_inf_B);
            is_inf_prod <= (is_inf_A | is_inf_B) & ~is_nan_prod;
            is_zero_prod <= (is_zero_A | is_zero_B) & ~is_nan_prod;
            
            if (is_nan_prod) begin
                // NaN result
                C <= {sign_prod, 5'b11111, 10'b1};
            end else if (is_inf_prod) begin
                // Infinity result
                C <= {sign_prod, 5'b11111, 10'b0};
            end else if (is_zero_prod) begin
                // Zero result
                C <= {sign_prod, 15'b0};
            end else begin
                // Normal case: apply rounding
                if (round_up) begin
                    frac_rounded <= frac_normalized[21:11] + 1;
                    exp_adj <= exp_normalized;
                    
                    // Handle carry from rounding
                    if (frac_rounded[10]) begin
                        frac_rounded <= {1'b0, frac_rounded[10:1]};
                        exp_adj <= exp_normalized + 1;
                    end
                end else begin
                    frac_rounded <= frac_normalized[21:11];
                    exp_adj <= exp_normalized;
                end
                
                // Check for overflow/underflow
                if (exp_adj >= EXP_MAX) begin
                    // Overflow to infinity
                    C <= {sign_prod, 5'b11111, 10'b0};
                end else if (exp_adj <= EXP_MIN) begin
                    // Underflow to zero
                    C <= {sign_prod, 15'b0};
                end else begin
                    // Normal result
                    C <= {sign_prod, exp_adj, frac_rounded[9:0]};
                end
            end
            
            Valid <= 1'b1;
        end else begin
            stage2_valid <= 1'b0;
            Valid <= 1'b0;
        end
    end
    
    // Input register and ready handling
    always @(posedge clk or negedge rst_n) begin
        if (!rst_n) begin
            A_reg <= 16'b0;
            B_reg <= 16'b0;
            ready_reg <= 1'b0;
            state <= 2'b0;
        end else begin
            case (state)
                2'b00: begin // Idle state
                    if (Ready) begin
                        A_reg <= A;
                        B_reg <= B;
                        ready_reg <= 1'b1;
                        state <= 2'b01;
                    end
                end
                2'b01: begin // Processing state
                    ready_reg <= 1'b0;
                    if (stage2_valid) begin
                        state <= 2'b00;
                    end
                end
                default: state <= 2'b00;
            endcase
        end
    end

endmodule
    
\end{lstlisting}


\newpage
\section{功能二：代码纠错}
将代码选中，右键选择Fix Code With Error Info的功能。
\begin{figure}[H]
    \centering
    \includegraphics[width=\textwidth]{代码修改.png}
\end{figure}
\begin{figure}[H]
    \centering
    \includegraphics[width=\textwidth]{代码修改结果.png}
    \caption{代码纠错}
\end{figure}

修改后的代码如下：
\begin{lstlisting}[frame=shadowbox,caption=vRTL251119111749\_415.v]
module fp_multiplier (
    input clk,
    input rst_n,
    input [15:0] A,
    input [15:0] B,
    input Ready,
    output reg [15:0] C,
    output reg Valid
);

    reg [15:0] computed_C;

    wire sign_A = A[15];
    wire [4:0] exp_A = A[14:10];
    wire [9:0] mant_A = A[9:0];

    wire sign_B = B[15];
    wire [4:0] exp_B = B[14:10];
    wire [9:0] mant_B = B[9:0];

    wire is_nan_A = (exp_A == 5'b11111) && (mant_A != 10'b0);
    wire is_inf_A = (exp_A == 5'b11111) && (mant_A == 10'b0);
    wire is_zero_A = (exp_A == 5'b0) && (mant_A == 10'b0);

    wire is_nan_B = (exp_B == 5'b11111) && (mant_B != 10'b0);
    wire is_inf_B = (exp_B == 5'b11111) && (mant_B == 10'b0);
    wire is_zero_B = (exp_B == 5'b0) && (mant_B == 10'b0);

    always @(*) begin
        computed_C = 16'b0;

        if (is_nan_A || is_nan_B) begin
            computed_C = {1'b0, 5'b11111, 10'b1};
        end else if ((is_inf_A && is_zero_B) || (is_zero_A && is_inf_B)) begin
            computed_C = {1'b0, 5'b11111, 10'b1};
        end else if (is_inf_A || is_inf_B) begin
            computed_C = {sign_A ^ sign_B, 5'b11111, 10'b0};
        end else if (is_zero_A || is_zero_B) begin
            computed_C = {sign_A ^ sign_B, 5'b0, 10'b0};
        end else begin
            wire [10:0] sig_A = (exp_A == 5'b0) ? {1'b0, mant_A} : {1'b1, mant_A};
            wire [10:0] sig_B = (exp_B == 5'b0) ? {1'b0, mant_B} : {1'b1, mant_B};
            wire [21:0] product_sig = sig_A * sig_B;

            wire [5:0] exp_A_adj = (exp_A == 5'b0) ? 6'b100001 : {1'b0, exp_A};
            wire [5:0] exp_B_adj = (exp_B == 5'b0) ? 6'b100001 : {1'b0, exp_B};
            wire [6:0] exp_sum = exp_A_adj + exp_B_adj;
            wire [6:0] exp_prod = exp_sum - 7'b0011110;

            reg [4:0] shift_amount;
            reg [21:0] normalized_sig;
            
            if (product_sig[21]) begin
                shift_amount = 5'b0;
                normalized_sig = product_sig;
            end else if (product_sig[20]) begin
                shift_amount = 5'b1;
                normalized_sig = product_sig << 1;
            end else if (product_sig[19]) begin
                shift_amount = 5'b10;
                normalized_sig = product_sig << 2;
            end else if (product_sig[18]) begin
                shift_amount = 5'b11;
                normalized_sig = product_sig << 3;
            end else begin
                shift_amount = 5'b0;
                normalized_sig = product_sig;
            end

            wire [6:0] exp_after_norm = exp_prod - {2'b0, shift_amount};

            wire [9:0] mant_part = normalized_sig[20:11];
            wire G = normalized_sig[10];
            wire R = normalized_sig[9];
            wire S = |normalized_sig[8:0];

            wire round_up;
            if (G && (R || S)) begin
                round_up = 1'b1;
            end else if (G && !(R || S) && mant_part[0]) begin
                round_up = 1'b1;
            end else begin
                round_up = 1'b0;
            end

            wire [10:0] rounded_mant = {1'b0, mant_part} + round_up;
            wire [9:0] final_mant;
            wire [6:0] final_exp;
            
            if (rounded_mant[10]) begin
                final_mant = rounded_mant[10:1];
                final_exp = exp_after_norm + 7'b1;
            end else begin
                final_mant = rounded_mant[9:0];
                final_exp = exp_after_norm;
            end

            if (final_exp[6] || final_exp[5:0] > 6'b001111) begin
                computed_C = {sign_A ^ sign_B, 5'b11111, 10'b0};
            end else if (final_exp[5:0] < 6'b000001) begin
                computed_C = {sign_A ^ sign_B, 5'b0, 10'b0};
            end else begin
                computed_C = {sign_A ^ sign_B, final_exp[4:0], final_mant};
            end
        end
    end

    always @(posedge clk or negedge rst_n) begin
        if (!rst_n) begin
            C <= 16'b0;
            Valid <= 1'b0;
        end else begin
            if (Ready) begin
                C <= computed_C;
                Valid <= 1'b1;
            end else begin
                Valid <= 1'b0;
            end
        end
    end
endmodule
\end{lstlisting}

对这个修改过的代码进行检查，发现还是有些地方出错，比如在always中使用wire。因此再对代码进行二次修改：
\begin{figure}[H]
    \centering
    \includegraphics[width=\textwidth]{代码修改结果2.png}
\end{figure}

第二次修改的结果如下：
\begin{lstlisting}[frame=shadowbox,caption=vRTL251119112055\_473.v]
module fp_multiplier (
    input clk,
    input rst_n,
    input [15:0] A,
    input [15:0] B,
    input Ready,
    output reg [15:0] C,
    output reg Valid
);

    reg [15:0] computed_C;

    wire sign_A = A[15];
    wire [4:0] exp_A = A[14:10];
    wire [9:0] mant_A = A[9:0];

    wire sign_B = B[15];
    wire [4:0] exp_B = B[14:10];
    wire [9:0] mant_B = B[9:0];

    wire is_nan_A = (exp_A == 5'b11111) && (mant_A != 10'b0);
    wire is_inf_A = (exp_A == 5'b11111) && (mant_A == 10'b0);
    wire is_zero_A = (exp_A == 5'b0) && (mant_A == 10'b0);

    wire is_nan_B = (exp_B == 5'b11111) && (mant_B != 10'b0);
    wire is_inf_B = (exp_B == 5'b11111) && (mant_B == 10'b0);
    wire is_zero_B = (exp_B == 5'b0) && (mant_B == 10'b0);

    // 将原来在always块内部的wire声明移到外部
    wire [10:0] sig_A = (exp_A == 5'b0) ? {1'b0, mant_A} : {1'b1, mant_A};
    wire [10:0] sig_B = (exp_B == 5'b0) ? {1'b0, mant_B} : {1'b1, mant_B};
    wire [21:0] product_sig = sig_A * sig_B;
    wire [5:0] exp_A_adj = (exp_A == 5'b0) ? 6'b100001 : {1'b0, exp_A};
    wire [5:0] exp_B_adj = (exp_B == 5'b0) ? 6'b100001 : {1'b0, exp_B};
    wire [6:0] exp_sum = exp_A_adj + exp_B_adj;
    wire [6:0] exp_prod = exp_sum - 7'b0011110;

    always @(*) begin
        computed_C = 16'b0;

        if (is_nan_A || is_nan_B) begin
            computed_C = {1'b0, 5'b11111, 10'b1};
        end else if ((is_inf_A && is_zero_B) || (is_zero_A && is_inf_B)) begin
            computed_C = {1'b0, 5'b11111, 10'b1};
        end else if (is_inf_A || is_inf_B) begin
            computed_C = {sign_A ^ sign_B, 5'b11111, 10'b0};
        end else if (is_zero_A || is_zero_B) begin
            computed_C = {sign_A ^ sign_B, 5'b0, 10'b0};
        end else begin
            reg [4:0] shift_amount;
            reg [21:0] normalized_sig;
            
            if (product_sig[21]) begin
                shift_amount = 5'b0;
                normalized_sig = product_sig;
            end else if (product_sig[20]) begin
                shift_amount = 5'b1;
                normalized_sig = product_sig << 1;
            end else if (product_sig[19]) begin
                shift_amount = 5'b10;
                normalized_sig = product_sig << 2;
            end else if (product_sig[18]) begin
                shift_amount = 5'b11;
                normalized_sig = product_sig << 3;
            end else begin
                shift_amount = 5'b0;
                normalized_sig = product_sig;
            end

            wire [6:0] exp_after_norm = exp_prod - {2'b0, shift_amount};

            wire [9:0] mant_part = normalized_sig[20:11];
            wire G = normalized_sig[10];
            wire R = normalized_sig[9];
            wire S = |normalized_sig[8:0];

            wire round_up;
            if (G && (R || S)) begin
                round_up = 1'b1;
            end else if (G && !(R || S) && mant_part[0]) begin
                round_up = 1'b1;
            end else begin
                round_up = 1'b0;
            end

            wire [10:0] rounded_mant = {1'b0, mant_part} + round_up;
            wire [9:0] final_mant;
            wire [6:0] final_exp;
            
            if (rounded_mant[10]) begin
                final_mant = rounded_mant[10:1];
                final_exp = exp_after_norm + 7'b1;
            end else begin
                final_mant = rounded_mant[9:0];
                final_exp = exp_after_norm;
            end

            if (final_exp[6] || final_exp[5:0] > 6'b001111) begin
                computed_C = {sign_A ^ sign_B, 5'b11111, 10'b0};
            end else if (final_exp[5:0] < 6'b000001) begin
                computed_C = {sign_A ^ sign_B, 5'b0, 10'b0};
            end else begin
                computed_C = {sign_A ^ sign_B, final_exp[4:0], final_mant};
            end
        end
    end

    always @(posedge clk or negedge rst_n) begin
        if (!rst_n) begin
            C <= 16'b0;
            Valid <= 1'b0;
        end else begin
            if (Ready) begin
                C <= computed_C;
                Valid <= 1'b1;
            end else begin
                Valid <= 1'b0;
            end
        end
    end
endmodule 
\end{lstlisting}


\newpage
\section{功能三：在聊天框进行代码纠错}
将代码右键Add to Chat，或者在Chat区添加文件。然后对模型提出修改任务和目标。这里我将任务一中补全的代码进行修改。
\begin{figure}[H]
    \centering
    \includegraphics[width=\textwidth]{聊天框纠错.png}
    \caption{聊天框纠错}
\end{figure}

将修改的结果插入到源代码中：

\begin{lstlisting}[frame=shadowbox]

//Revised version

module fp_multiplier (
    input clk,
    input rst_n,
    input [15:0] A,
    input [15:0] B,
    input Ready,
    output reg [15:0] C,
    output reg Valid
);

    // Internal registers for pipeline stages
    reg [15:0] A_reg, B_reg;
    reg ready_reg;
    reg [1:0] state;
    
    // Floating point components extraction
    wire sign_A = A_reg[15];
    wire [4:0] exp_A = A_reg[14:10];
    wire [9:0] frac_A = A_reg[9:0];
    
    wire sign_B = B_reg[15];
    wire [4:0] exp_B = B_reg[14:10];
    wire [9:0] frac_B = B_reg[9:0];
    
    // Stage 1 signals
    reg stage1_valid;
    reg sign_prod;
    reg [5:0] exp_sum;
    reg [21:0] frac_prod;
    reg is_zero_A, is_zero_B;
    reg is_inf_A, is_inf_B;
    reg is_nan_A, is_nan_B;
    
    // Stage 2 signals
    reg stage2_valid;
    reg [5:0] exp_normalized;
    reg [21:0] frac_normalized;
    reg guard_bit, round_bit, sticky_bit;
    reg round_up;
    reg is_zero_prod, is_inf_prod, is_nan_prod;
    
    // Constants
    localparam [4:0] EXP_BIAS = 5'd15;
    localparam [4:0] EXP_MAX = 5'd31;
    localparam [4:0] EXP_MIN = 5'd0;
    
    // Special case detection (combinational)
    always @(*) begin
        is_zero_A = (exp_A == 5'b0) && (frac_A == 10'b0);
        is_inf_A = (exp_A == 5'b11111) && (frac_A == 10'b0);
        is_nan_A = (exp_A == 5'b11111) && (frac_A != 10'b0);
        
        is_zero_B = (exp_B == 5'b0) && (frac_B == 10'b0);
        is_inf_B = (exp_B == 5'b11111) && (frac_B == 10'b0);
        is_nan_B = (exp_B == 5'b11111) && (frac_B != 10'b0);
    end
    
    // Pipeline stage 1: Multiplication and special case handling
    always @(posedge clk or negedge rst_n) begin
        if (!rst_n) begin
            stage1_valid <= 1'b0;
            sign_prod <= 1'b0;
            exp_sum <= 6'b0;
            frac_prod <= 22'b0;
            exp_normalized <= 6'b0;
            frac_normalized <= 22'b0;
            guard_bit <= 1'b0;
            round_bit <= 1'b0;
            sticky_bit <= 1'b0;
            round_up <= 1'b0;
        end else if (ready_reg) begin
            stage1_valid <= 1'b1;
            
            // Sign calculation
            sign_prod <= sign_A ^ sign_B;
            
            // Exponent calculation (subtract bias)
            exp_sum <= {1'b0, exp_A} + {1'b0, exp_B} - EXP_BIAS;
            
            // Fraction multiplication (with implicit 1)
            if (is_zero_A || is_zero_B) begin
                frac_prod <= 22'b0;
            end else begin
                frac_prod <= {1'b1, frac_A} * {1'b1, frac_B};
            end
            
            // Normalization logic (moved from combinational to sequential)
            if (frac_prod[21]) begin
                frac_normalized <= frac_prod >> 1;
                exp_normalized <= exp_sum + 1;
            end else begin
                frac_normalized <= frac_prod;
                exp_normalized <= exp_sum;
            end
            
            // Rounding bits calculation
            guard_bit <= frac_normalized[10];
            round_bit <= frac_normalized[9];
            sticky_bit <= |frac_normalized[8:0];
            
            // Round to nearest even
            round_up <= guard_bit & (round_bit | sticky_bit | frac_normalized[11]);
            
        end else begin
            stage1_valid <= 1'b0;
        end
    end
    
    // Pipeline stage 2: Rounding and final result assembly
    always @(posedge clk or negedge rst_n) begin
        if (!rst_n) begin
            stage2_valid <= 1'b0;
            C <= 16'b0;
            Valid <= 1'b0;
            is_zero_prod <= 1'b0;
            is_inf_prod <= 1'b0;
            is_nan_prod <= 1'b0;
        end else if (stage1_valid) begin
            stage2_valid <= 1'b1;
            
            // Handle special cases
            is_nan_prod <= is_nan_A | is_nan_B | (is_inf_A & is_zero_B) | (is_zero_A & is_inf_B);
            is_inf_prod <= (is_inf_A | is_inf_B) & ~is_nan_prod;
            is_zero_prod <= (is_zero_A | is_zero_B) & ~is_nan_prod;
            
            if (is_nan_prod) begin
                // NaN result
                C <= {sign_prod, 5'b11111, 10'b1};
            end else if (is_inf_prod) begin
                // Infinity result
                C <= {sign_prod, 5'b11111, 10'b0};
            end else if (is_zero_prod) begin
                // Zero result
                C <= {sign_prod, 15'b0};
            end else begin
                // Normal case: apply rounding with carry handling
                reg [10:0] temp_frac_rounded;
                reg [5:0] temp_exp_adj;
                
                if (round_up) begin
                    temp_frac_rounded = frac_normalized[21:11] + 1;
                    temp_exp_adj = exp_normalized;
                    
                    // Handle carry from rounding
                    if (temp_frac_rounded[10]) begin
                        temp_frac_rounded = {1'b0, temp_frac_rounded[10:1]};
                        temp_exp_adj = temp_exp_adj + 1;
                    end
                end else begin
                    temp_frac_rounded = frac_normalized[21:11];
                    temp_exp_adj = exp_normalized;
                end
                
                // Check for overflow/underflow
                if (temp_exp_adj >= EXP_MAX) begin
                    // Overflow to infinity
                    C <= {sign_prod, 5'b11111, 10'b0};
                end else if (temp_exp_adj <= EXP_MIN) begin
                    // Underflow to zero
                    C <= {sign_prod, 15'b0};
                end else begin
                    // Normal result
                    C <= {sign_prod, temp_exp_adj[4:0], temp_frac_rounded[9:0]};
                end
            end
            
            Valid <= 1'b1;
        end else begin
            stage2_valid <= 1'b0;
            Valid <= 1'b0;
        end
    end
    
    // Input register and state machine
    always @(posedge clk or negedge rst_n) begin
        if (!rst_n) begin
            A_reg <= 16'b0;
            B_reg <= 16'b0;
            ready_reg <= 1'b0;
            state <= 2'b00;
        end else begin
            case (state)
                2'b00: begin // Idle state
                    ready_reg <= 1'b0;
                    if (Ready) begin
                        A_reg <= A;
                        B_reg <= B;
                        ready_reg <= 1'b1;
                        state <= 2'b01;
                    end
                end
                2'b01: begin // Stage 1 processing
                    ready_reg <= 1'b0;
                    state <= 2'b10;
                end
                2'b10: begin // Stage 2 processing
                    if (stage2_valid) begin
                        state <= 2'b00;
                    end
                end
                default: begin
                    state <= 2'b00;
                    ready_reg <= 1'b0;
                end
            endcase
        end
    end

endmodule
\end{lstlisting}


\section{功能四：基于Specification的RTL生成}
因为我想看看该AI-EDA工具的生成能力，因此这里采用的是助教在HW1中提供的Zero Shot Prompt3，让其有更多
可以自由发挥的设计空间。
\begin{lstlisting}[frame=shadowbox]
Please help me generate a floating-point multiplier unit, including RTL code and TestBench. 

### Problem Description

ModelName: fp_multiplier

InputNode:

- clk: Clock signal (rising edge effective).
- rst_n: Active-low reset signal.
- A: 16-bit FP16 format input, operand one.
- B: 16-bit FP16 format input, operand two.
- Ready: Input data ready signal, high effective.

OutputNode:

- C: 16-bit FP16 format output.
- Valid: Device idle and result ready signal, high effective.

Specification:

- Support IEEE 754 Half-Precision Floating-Point Format (FP16):
  - Handle all special cases (zero, infinity, NaN) as specified by IEEE 754.
  - Ensure correct rounding to the nearest even value.
- Use valid and ready signals for handshake. When valid is high, the input operands A and B are ready. When ready is high, the result C is valid.
- TestBench should contain the following test conditions:
  - Basic Multiplication: Positive numbers (e.g., 2.0 * 3.0); Negative numbers (e.g., -2.0 * 3.0); Mixed signs (e.g., -2.0 * -3.0);
  - Zero Multiplication: Multiplication involving zero (e.g., 0.0 * 2.0);
  - Infinity Multiplication: Multiplication involving infinity (e.g., inf * 2.0);
  - NaN Multiplication: Multiplication involving NaN (e.g., NaN * 2.0);
  - Subnormal Numbers: Multiplication involving very small numbers (e.g., 5.96e-8 * 5.96e-8);
  - Overflow and Underflow: Test cases that result in overflow or underflow;
  - Edge Cases: Multiplication of maximum representable values (e.g., 65504.0 * 1.5); Multiplication of minimum representable values (e.g., 5.96e-8 * 5.96e-8);


### RTL

```verilog
module fp_multiplier
(
    input wire clk,
    input wire rst_n,
    input wire [15:0] A,
    input wire [15:0] B,
    output wire [15:0] C,
    input wire valid,
    output wire ready
)
```

### TestBench

```systemverilog
module fp_multiplier_tb;
```


### Tips

Here is some tips that you need to follow:

1. When responsing, please fill in RTL and TestBench respectively in the code box starting with "### RTL" and "### TestBench".
2. Use Verilog for RTL and SystemVerilog for TestBench.
3. We use the -g2005sv specification supported by iverilog for compilation. Please make sure that the verilog used can be supported.
4. Unpacked structs not supported in iverilog, do not use typedef struct.
5. When using the $bitstoreal function, please make sure that the parameter is 64bit.
6. Save the simulated waveform diagram with $dumpfile and $dumpvars instructions. Save the vcd file as Out/wave.vcd
7. In Testbench, please verify whether the function is correct. If you pass all the tests, please output TestSuccessed, otherwise you will output TestFailed.


\end{lstlisting}

将这个prompt放入Chat聊天框中，得到的生成结果如下。
\begin{figure}[H]
    \centering
    \includegraphics[width=\textwidth]{spec生成rtl.png}
\end{figure}

\begin{lstlisting}[frame=shadowbox,caption=fp\_multiplier\_generated.v]
module fp_multiplier
(
    input wire clk,
    input wire rst_n,
    input wire [15:0] A,
    input wire [15:0] B,
    output reg [15:0] C,
    input wire valid,
    output reg ready
);

    // Internal registers for pipeline
    reg [15:0] A_reg, B_reg;
    reg valid_reg;
    reg [1:0] state;
    
    // FP16 format: 1-bit sign, 5-bit exponent, 10-bit fraction
    wire sign_a = A[15];
    wire sign_b = B[15];
    wire [4:0] exp_a = A[14:10];
    wire [4:0] exp_b = B[14:10];
    wire [9:0] frac_a = A[9:0];
    wire [9:0] frac_b = B[9:0];
    
    // Special case detection
    wire a_is_zero = (exp_a == 5'b0) && (frac_a == 10'b0);
    wire b_is_zero = (exp_b == 5'b0) && (frac_b == 10'b0);
    wire a_is_inf = (exp_a == 5'b11111) && (frac_a == 10'b0);
    wire b_is_inf = (exp_b == 5'b11111) && (frac_b == 10'b0);
    wire a_is_nan = (exp_a == 5'b11111) && (frac_a != 10'b0);
    wire b_is_nan = (exp_b == 5'b11111) && (frac_b != 10'b0);
    wire a_is_subnormal = (exp_a == 5'b0) && (frac_a != 10'b0);
    wire b_is_subnormal = (exp_b == 5'b0) && (frac_b != 10'b0);
    
    // Pipeline stage 1: Special case handling and mantissa multiplication preparation
    reg stage1_valid;
    reg stage1_sign;
    reg [5:0] stage1_exp_sum; // 6 bits for sum + bias adjustment
    reg [20:0] stage1_mant_product; // 11x11 = 22 bits, but we use 21 for calculation
    reg stage1_is_zero, stage1_is_inf, stage1_is_nan;
    
    // Pipeline stage 2: Normalization and rounding
    reg stage2_valid;
    reg [15:0] stage2_result;
    
    // FSM states
    localparam IDLE = 2'b00;
    localparam STAGE1 = 2'b01;
    localparam STAGE2 = 2'b10;
    
    always @(posedge clk or negedge rst_n) begin
        if (!rst_n) begin
            state <= IDLE;
            ready <= 1'b1;
            valid_reg <= 1'b0;
            A_reg <= 16'b0;
            B_reg <= 16'b0;
            stage1_valid <= 1'b0;
            stage2_valid <= 1'b0;
            C <= 16'b0;
        end else begin
            case (state)
                IDLE: begin
                    ready <= 1'b1;
                    if (valid && ready) begin
                        A_reg <= A;
                        B_reg <= B;
                        valid_reg <= 1'b1;
                        state <= STAGE1;
                        ready <= 1'b0;
                        
                        // Stage 1 calculations
                        stage1_sign <= sign_a ^ sign_b;
                        
                        // Handle special cases
                        if (a_is_nan || b_is_nan) begin
                            stage1_is_nan <= 1'b1;
                            stage1_is_zero <= 1'b0;
                            stage1_is_inf <= 1'b0;
                            stage1_exp_sum <= 6'b0;
                            stage1_mant_product <= 21'b0;
                        end else if ((a_is_zero && b_is_inf) || (a_is_inf && b_is_zero)) begin
                            stage1_is_nan <= 1'b1;
                            stage1_is_zero <= 1'b0;
                            stage1_is_inf <= 1'b0;
                            stage1_exp_sum <= 6'b0;
                            stage1_mant_product <= 21'b0;
                        end else if (a_is_zero || b_is_zero) begin
                            stage1_is_zero <= 1'b1;
                            stage1_is_inf <= 1'b0;
                            stage1_is_nan <= 1'b0;
                            stage1_exp_sum <= 6'b0;
                            stage1_mant_product <= 21'b0;
                        end else if (a_is_inf || b_is_inf) begin
                            stage1_is_inf <= 1'b1;
                            stage1_is_zero <= 1'b0;
                            stage1_is_nan <= 1'b0;
                            stage1_exp_sum <= 6'b0;
                            stage1_mant_product <= 21'b0;
                        end else begin
                            // Normal multiplication
                            stage1_is_zero <= 1'b0;
                            stage1_is_inf <= 1'b0;
                            stage1_is_nan <= 1'b0;
                            
                            // Add exponents and subtract bias (15 for FP16)
                            stage1_exp_sum <= {1'b0, exp_a} + {1'b0, exp_b} - 6'd15;
                            
                            // Multiply mantissas (with hidden bit)
                            wire [10:0] mant_a = (a_is_subnormal) ? {1'b0, frac_a} : {1'b1, frac_a};
                            wire [10:0] mant_b = (b_is_subnormal) ? {1'b0, frac_b} : {1'b1, frac_b};
                            stage1_mant_product <= mant_a * mant_b;
                        end
                        
                        stage1_valid <= 1'b1;
                    end
                end
                
                STAGE1: begin
                    if (stage1_valid) begin
                        // Stage 2: Normalization and rounding
                        if (stage1_is_nan) begin
                            stage2_result <= {stage1_sign, 5'b11111, 10'b1}; // Quiet NaN
                        end else if (stage1_is_inf) begin
                            stage2_result <= {stage1_sign, 5'b11111, 10'b0}; // Infinity
                        end else if (stage1_is_zero) begin
                            stage2_result <= {stage1_sign, 15'b0}; // Zero
                        end else begin
                            // Normal case: normalize and round
                            reg [4:0] final_exp;
                            reg [9:0] final_frac;
                            reg [20:0] normalized_mant;
                            
                            // Check for overflow in exponent
                            if (stage1_exp_sum[5] || (stage1_exp_sum[4:0] > 5'd30)) begin
                                // Overflow - return infinity
                                stage2_result <= {stage1_sign, 5'b11111, 10'b0};
                            end else if (stage1_exp_sum[4:0] == 5'b0) begin
                                // Underflow - return zero
                                stage2_result <= {stage1_sign, 15'b0};
                            end else begin
                                // Normal case
                                normalized_mant = stage1_mant_product;
                                final_exp = stage1_exp_sum[4:0];
                                
                                // Normalize: product is in [1,4) range, so we may need to shift
                                if (normalized_mant[20]) begin
                                    normalized_mant = normalized_mant >> 1;
                                    final_exp = final_exp + 1;
                                    
                                    // Check for exponent overflow after normalization
                                    if (final_exp > 5'd30) begin
                                        stage2_result <= {stage1_sign, 5'b11111, 10'b0};
                                    end else begin
                                        // Round to nearest even (simplified)
                                        final_frac = normalized_mant[19:10];
                                        if (normalized_mant[9] && (normalized_mant[8:0] != 9'b0 || normalized_mant[10])) begin
                                            final_frac = final_frac + 1;
                                            if (final_frac == 10'b0) begin // Carry occurred
                                                final_exp = final_exp + 1;
                                                if (final_exp > 5'd30) begin
                                                    stage2_result <= {stage1_sign, 5'b11111, 10'b0};
                                                end else begin
                                                    stage2_result <= {stage1_sign, final_exp, final_frac};
                                                end
                                            end else begin
                                                stage2_result <= {stage1_sign, final_exp, final_frac};
                                            end
                                        end else begin
                                            stage2_result <= {stage1_sign, final_exp, final_frac};
                                        end
                                    end
                                end else begin
                                    // No normalization needed
                                    final_frac = normalized_mant[18:9];
                                    if (normalized_mant[8] && (normalized_mant[7:0] != 8'b0 || normalized_mant[9])) begin
                                        final_frac = final_frac + 1;
                                        if (final_frac == 10'b0) begin // Carry occurred
                                            final_exp = final_exp + 1;
                                            if (final_exp > 5'd30) begin
                                                stage2_result <= {stage1_sign, 5'b11111, 10'b0};
                                            end else begin
                                                stage2_result <= {stage1_sign, final_exp, final_frac};
                                            end
                                        end else begin
                                            stage2_result <= {stage1_sign, final_exp, final_frac};
                                        end
                                    end else begin
                                        stage2_result <= {stage1_sign, final_exp, final_frac};
                                    end
                                end
                            end
                        end
                        
                        stage2_valid <= 1'b1;
                        state <= STAGE2;
                    end
                end
                
                STAGE2: begin
                    if (stage2_valid) begin
                        C <= stage2_result;
                        stage2_valid <= 1'b0;
                        stage1_valid <= 1'b0;
                        valid_reg <= 1'b0;
                        state <= IDLE;
                    end
                end
            endcase
        end
    end

endmodule
\end{lstlisting}

\section{功能五：TestBench生成及仿真结果}

\begin{lstlisting}[frame=shadowbox,caption=tb\_fp\_multiplier\_generated.sv]
module fp_multiplier_tb;
    reg clk;
    reg rst_n;
    reg [15:0] A;
    reg [15:0] B;
    wire [15:0] C;
    reg valid;
    wire ready;
    
    // Test results
    integer test_count = 0;
    integer pass_count = 0;
    integer fail_count = 0;
    
    // Instantiate DUT
    fp_multiplier dut (
        .clk(clk),
        .rst_n(rst_n),
        .A(A),
        .B(B),
        .C(C),
        .valid(valid),
        .ready(ready)
    );
    
    // Clock generation
    always #5 clk = ~clk;
    
    // FP16 value definitions
    parameter FP16_ZERO = 16'h0000;
    parameter FP16_ONE = 16'h3C00;
    parameter FP16_TWO = 16'h4000;
    parameter FP16_THREE = 16'h4200;
    parameter FP16_NEG_TWO = 16'hC000;
    parameter FP16_NEG_THREE = 16'hC200;
    parameter FP16_INF = 16'h7C00;
    parameter FP16_NEG_INF = 16'hFC00;
    parameter FP16_NAN = 16'h7E00;
    parameter FP16_MAX = 16'h7BFF;  // 65504
    parameter FP16_MIN = 16'h0001;  // ~5.96e-8
    
    // Function to convert FP16 to real (simplified)
    function real fp16_to_real;
        input [15:0] fp16;
        reg sign;
        reg [4:0] exp;
        reg [9:0] frac;
        real result;
        integer i;
        real fraction_value;
    begin
        sign = fp16[15];
        exp = fp16[14:10];
        frac = fp16[9:0];
        
        if (exp == 5'b11111) begin
            if (frac == 10'b0) begin
                fp16_to_real = (sign) ? -$bitstoreal(64'h7FF0000000000000) : $bitstoreal(64'h7FF0000000000000);
            end else begin
                fp16_to_real = $bitstoreal(64'h7FF8000000000000); // NaN
            end
        end else if (exp == 5'b0) then begin
            // Subnormal
            fraction_value = 0.0;
            for (i = 0; i < 10; i = i + 1) begin
                if (frac[9-i]) fraction_value = fraction_value + (2.0 ** (-i-1));
            end
            result = fraction_value * (2.0 ** -14);
            fp16_to_real = (sign) ? -result : result;
        end else begin
            // Normal
            fraction_value = 1.0;
            for (i = 0; i < 10; i = i + 1) begin
                if (frac[9-i]) fraction_value = fraction_value + (2.0 ** (-i-1));
            end
            result = fraction_value * (2.0 ** (exp - 15));
            fp16_to_real = (sign) ? -result : result;
        end
    end
    endfunction
    
    // Test task
    task test_multiply;
        input [15:0] a_val;
        input [15:0] b_val;
        input [15:0] expected;
        string test_name;
        real a_real, b_real, expected_real, result_real;
        reg [63:0] expected_bits, result_bits;
    begin
        test_count = test_count + 1;
        
        // Wait for ready
        @(posedge clk);
        while (!ready) @(posedge clk);
        
        // Apply inputs
        A = a_val;
        B = b_val;
        valid = 1'b1;
        
        // Wait for computation
        @(posedge clk);
        valid = 1'b0;
        
        // Wait for result (3 cycles for pipeline)
        repeat(3) @(posedge clk);
        
        // Convert to real for comparison (handle special cases)
        a_real = fp16_to_real(a_val);
        b_real = fp16_to_real(b_val);
        expected_real = fp16_to_real(expected);
        result_real = fp16_to_real(C);
        
        expected_bits = $realtobits(expected_real);
        result_bits = $realtobits(result_real);
        
        // Check result
        if (expected_bits === result_bits) begin
            pass_count = pass_count + 1;
            $display("PASS: Test %0d - A=%h, B=%h, Got=%h, Expected=%h", 
                     test_count, a_val, b_val, C, expected);
        end else begin
            fail_count = fail_count + 1;
            $display("FAIL: Test %0d - A=%h, B=%h, Got=%h, Expected=%h", 
                     test_count, a_val, b_val, C, expected);
            $display("      A_real=%f, B_real=%f, Got_real=%f, Expected_real=%f",
                     a_real, b_real, result_real, expected_real);
        end
    end
    endtask
    
    initial begin
        // Initialize
        clk = 0;
        rst_n = 0;
        valid = 0;
        A = 0;
        B = 0;
        
        // Generate waveform
        $dumpfile("Out/wave.vcd");
        $dumpvars(0, fp_multiplier_tb);
        
        // Reset
        #10 rst_n = 1;
        
        $display("Starting FP16 Multiplier Tests...");
        
        // Test 1: Basic Multiplication - Positive numbers
        test_multiply(FP16_TWO, FP16_THREE, 16'h4400); // 2.0 * 3.0 = 6.0
        
        // Test 2: Negative numbers  
        test_multiply(FP16_NEG_TWO, FP16_THREE, 16'hC400); // -2.0 * 3.0 = -6.0
        
        // Test 3: Mixed signs
        test_multiply(FP16_NEG_TWO, FP16_NEG_THREE, 16'h4400); // -2.0 * -3.0 = 6.0
        
        // Test 4: Zero multiplication
        test_multiply(FP16_ZERO, FP16_TWO, FP16_ZERO); // 0.0 * 2.0 = 0.0
        test_multiply(FP16_TWO, FP16_ZERO, FP16_ZERO); // 2.0 * 0.0 = 0.0
        
        // Test 5: Infinity multiplication
        test_multiply(FP16_INF, FP16_TWO, FP16_INF); // inf * 2.0 = inf
        test_multiply(FP16_NEG_INF, FP16_THREE, FP16_NEG_INF); // -inf * 3.0 = -inf
        
        // Test 6: NaN multiplication
        test_multiply(FP16_NAN, FP16_TWO, FP16_NAN); // NaN * 2.0 = NaN
        
        // Test 7: Infinity * Zero = NaN
        test_multiply(FP16_INF, FP16_ZERO, FP16_NAN);
        test_multiply(FP16_ZERO, FP16_INF, FP16_NAN);
        
        // Test 8: Subnormal numbers
        test_multiply(FP16_MIN, FP16_MIN, FP16_ZERO); // Very small * very small = 0 (underflow)
        
        // Test 9: Overflow
        test_multiply(FP16_MAX, FP16_TWO, FP16_INF); // 65504 * 2.0 = inf (overflow)
        
        // Test 10: Edge case - maximum representable
        test_multiply(FP16_MAX, FP16_ONE, FP16_MAX); // 65504 * 1.0 = 65504
        
        // Final summary
        $display("\nTest Summary:");
        $display("Total Tests: %0d", test_count);
        $display("Passed: %0d", pass_count);
        $display("Failed: %0d", fail_count);
        
        if (fail_count == 0) begin
            $display("TestSuccessed");
        end else begin
            $display("TestFailed");
        end
        
        $finish;
    end

endmodule
\end{lstlisting}

在FILES中将这两个文件添加进来，运行仿真。
\begin{figure}[H]
    \centering
    \includegraphics[width=\textwidth]{仿真结果.png}
    \caption{第一次仿真结果}
\end{figure}

可以看到，这一版的代码可以运行，但是很多测试点的结果都不对。因此我们在聊天框将这个仿真结果输入给AI，并让其
帮我们修改代码。
\begin{figure}[H]
    \centering
    \includegraphics[width=\textwidth]{结合仿真结果修改生成代码.png}
    \caption{结合仿真结果修改生成代码}
\end{figure}

生成的代码如下：
\begin{lstlisting}[frame=shadowbox,caption=fp\_multiplier\_generated3.sv]
module fp_multiplier
(
    input wire clk,
    input wire rst_n,
    input wire [15:0] A,
    input wire [15:0] B,
    output reg [15:0] C,
    input wire valid,
    output reg ready
);

    // Internal registers for pipeline
    reg [15:0] A_reg, B_reg;
    reg valid_reg;
    reg [1:0] state;
    
    // FP16 format: 1-bit sign, 5-bit exponent, 10-bit fraction
    wire sign_a = A[15];
    wire sign_b = B[15];
    wire [4:0] exp_a = A[14:10];
    wire [4:0] exp_b = B[14:10];
    wire [9:0] frac_a = A[9:0];
    wire [9:0] frac_b = B[9:0];
    
    // Special case detection
    wire a_is_zero = (exp_a == 5'b0) && (frac_a == 10'b0);
    wire b_is_zero = (exp_b == 5'b0) && (frac_b == 10'b0);
    wire a_is_inf = (exp_a == 5'b11111) && (frac_a == 10'b0);
    wire b_is_inf = (exp_b == 5'b11111) && (frac_b == 10'b0);
    wire a_is_nan = (exp_a == 5'b11111) && (frac_a != 10'b0);
    wire b_is_nan = (exp_b == 5'b11111) && (frac_b != 10'b0);
    wire a_is_subnormal = (exp_a == 5'b0) && (frac_a != 10'b0);
    wire b_is_subnormal = (exp_b == 5'b0) && (frac_b != 10'b0);
    
    // Mantissa calculation wires
    wire [10:0] mant_a = (a_is_subnormal || a_is_zero) ? {1'b0, frac_a} : {1'b1, frac_a};
    wire [10:0] mant_b = (b_is_subnormal || b_is_zero) ? {1'b0, frac_b} : {1'b1, frac_b};
    
    // Pipeline stage 1: Special case handling and mantissa multiplication preparation
    reg stage1_valid;
    reg stage1_sign;
    reg [5:0] stage1_exp_sum; // 6 bits for sum + bias adjustment
    reg [21:0] stage1_mant_product; // 11x11 = 22 bits
    reg stage1_is_zero, stage1_is_inf, stage1_is_nan;
    
    // Pipeline stage 2: Normalization and rounding
    reg stage2_valid;
    reg [15:0] stage2_result;
    
    // FSM states
    localparam IDLE = 2'b00;
    localparam STAGE1 = 2'b01;
    localparam STAGE2 = 2'b10;
    
    always @(posedge clk or negedge rst_n) begin
        if (!rst_n) begin
            state <= IDLE;
            ready <= 1'b1;
            valid_reg <= 1'b0;
            A_reg <= 16'b0;
            B_reg <= 16'b0;
            stage1_valid <= 1'b0;
            stage2_valid <= 1'b0;
            C <= 16'b0;
        end else begin
            case (state)
                IDLE: begin
                    ready <= 1'b1;
                    if (valid && ready) begin
                        A_reg <= A;
                        B_reg <= B;
                        valid_reg <= 1'b1;
                        state <= STAGE1;
                        ready <= 1'b0;
                        
                        // Stage 1 calculations
                        stage1_sign <= sign_a ^ sign_b;
                        
                        // Handle special cases
                        if (a_is_nan || b_is_nan) begin
                            stage1_is_nan <= 1'b1;
                            stage1_is_zero <= 1'b0;
                            stage1_is_inf <= 1'b0;
                            stage1_exp_sum <= 6'b0;
                            stage1_mant_product <= 22'b0;
                        end else if ((a_is_zero && b_is_inf) || (a_is_inf && b_is_zero)) begin
                            stage1_is_nan <= 1'b1;
                            stage1_is_zero <= 1'b0;
                            stage1_is_inf <= 1'b0;
                            stage1_exp_sum <= 6'b0;
                            stage1_mant_product <= 22'b0;
                        end else if (a_is_zero || b_is_zero) begin
                            stage1_is_zero <= 1'b1;
                            stage1_is_inf <= 1'b0;
                            stage1_is_nan <= 1'b0;
                            stage1_exp_sum <= 6'b0;
                            stage1_mant_product <= 22'b0;
                        end else if (a_is_inf || b_is_inf) begin
                            stage1_is_inf <= 1'b1;
                            stage1_is_zero <= 1'b0;
                            stage1_is_nan <= 1'b0;
                            stage1_exp_sum <= 6'b0;
                            stage1_mant_product <= 22'b0;
                        end else begin
                            // Normal multiplication
                            stage1_is_zero <= 1'b0;
                            stage1_is_inf <= 1'b0;
                            stage1_is_nan <= 1'b0;
                            
                            // Add exponents and subtract bias (15 for FP16)
                            // For normal numbers: exp = stored_exp - 15
                            // So product_exp = (exp_a - 15) + (exp_b - 15) + 15 = exp_a + exp_b - 15
                            stage1_exp_sum <= {1'b0, exp_a} + {1'b0, exp_b} - 6'd15;
                            
                            // Multiply mantissas (with hidden bit)
                            stage1_mant_product <= mant_a * mant_b;
                        end
                        
                        stage1_valid <= 1'b1;
                    end
                end
                
                STAGE1: begin
                    if (stage1_valid) begin
                        // Stage 2: Normalization and rounding
                        if (stage1_is_nan) begin
                            // Return canonical NaN
                            stage2_result <= {stage1_sign, 5'b11111, 10'b1000000000};
                        end else if (stage1_is_inf) begin
                            stage2_result <= {stage1_sign, 5'b11111, 10'b0}; // Infinity
                        end else if (stage1_is_zero) begin
                            stage2_result <= {stage1_sign, 15'b0}; // Zero
                        end else begin
                            // Normal case: normalize and round
                            reg [5:0] temp_exp;
                            reg [9:0] temp_frac;
                            reg [21:0] temp_mant;
                            
                            temp_exp = stage1_exp_sum;
                            temp_mant = stage1_mant_product;
                            
                            // Normalize: product is in [1,4) range, so we may need to shift
                            if (temp_mant[21]) begin
                                // Product >= 2, need to shift right
                                temp_mant = temp_mant >> 1;
                                temp_exp = temp_exp + 1;
                            end
                            
                            // Check for overflow/underflow
                            if (temp_exp[5] || temp_exp[4:0] > 5'd30) begin
                                // Overflow - return infinity
                                stage2_result <= {stage1_sign, 5'b11111, 10'b0};
                            end else if (temp_exp[4:0] == 5'b0) begin
                                // Underflow - return zero
                                stage2_result <= {stage1_sign, 15'b0};
                            end else begin
                                // Extract fraction (ignore the hidden bit at position 20)
                                temp_frac = temp_mant[20:11];
                                
                                // Round to nearest even
                                if (temp_mant[10] && (temp_mant[9:0] != 10'b0 || temp_mant[11])) begin
                                    temp_frac = temp_frac + 1;
                                    if (temp_frac == 10'b0) begin // Carry to exponent
                                        temp_exp = temp_exp + 1;
                                        if (temp_exp > 5'd30) begin
                                            stage2_result <= {stage1_sign, 5'b11111, 10'b0};
                                        end else begin
                                            stage2_result <= {stage1_sign, temp_exp[4:0], temp_frac};
                                        end
                                    end else begin
                                        stage2_result <= {stage1_sign, temp_exp[4:0], temp_frac};
                                    end
                                end else begin
                                    stage2_result <= {stage1_sign, temp_exp[4:0], temp_frac};
                                end
                            end
                        end
                        
                        stage2_valid <= 1'b1;
                        state <= STAGE2;
                    end
                end
                
                STAGE2: begin
                    if (stage2_valid) begin
                        C <= stage2_result;
                        stage2_valid <= 1'b0;
                        stage1_valid <= 1'b0;
                        valid_reg <= 1'b0;
                        state <= IDLE;
                    end
                end
            endcase
        end
    end

endmodule
\end{lstlisting}
\begin{lstlisting}[frame=shadowbox,caption=vRTL251119114743\_749.v]
odule fp_multiplier_tb;
    reg clk;
    reg rst_n;
    reg [15:0] A;
    reg [15:0] B;
    wire [15:0] C;
    reg valid;
    wire ready;
    
    // Test results
    integer test_count = 0;
    integer pass_count = 0;
    integer fail_count = 0;
    
    // Instantiate DUT
    fp_multiplier dut (
        .clk(clk),
        .rst_n(rst_n),
        .A(A),
        .B(B),
        .C(C),
        .valid(valid),
        .ready(ready)
    );
    
    // Clock generation
    always #5 clk = ~clk;
    
    // FP16 value definitions
    parameter FP16_ZERO = 16'h0000;
    parameter FP16_ONE = 16'h3C00;
    parameter FP16_TWO = 16'h4000;
    parameter FP16_THREE = 16'h4200;
    parameter FP16_NEG_TWO = 16'hC000;
    parameter FP16_NEG_THREE = 16'hC200;
    parameter FP16_INF = 16'h7C00;
    parameter FP16_NEG_INF = 16'hFC00;
    parameter FP16_NAN = 16'h7E00;
    parameter FP16_MAX = 16'h7BFF;  // 65504
    parameter FP16_MIN = 16'h0001;  // ~5.96e-8
    parameter FP16_SIX = 16'h4600;  // 6.0
    
    // Function to convert FP16 to real (simplified)
    function real fp16_to_real;
        input [15:0] fp16;
        reg sign;
        reg [4:0] exp;
        reg [9:0] frac;
        real result;
        integer i;
        real fraction_value;
    begin
        sign = fp16[15];
        exp = fp16[14:10];
        frac = fp16[9:0];
        
        if (exp == 5'b11111) begin
            if (frac == 10'b0) begin
                fp16_to_real = (sign) ? -$bitstoreal(64'h7FF0000000000000) : $bitstoreal(64'h7FF0000000000000);
            end else begin
                fp16_to_real = $bitstoreal(64'h7FF8000000000000); // NaN
            end
        end else if (exp == 5'b0) begin
            // Subnormal
            fraction_value = 0.0;
            for (i = 0; i < 10; i = i + 1) begin
                if (frac[9-i]) fraction_value = fraction_value + (2.0 ** (-i-1));
            end
            result = fraction_value * (2.0 ** -14);
            fp16_to_real = (sign) ? -result : result;
        end else begin
            // Normal
            fraction_value = 1.0;
            for (i = 0; i < 10; i = i + 1) begin
                if (frac[9-i]) fraction_value = fraction_value + (2.0 ** (-i-1));
            end
            result = fraction_value * (2.0 ** (exp - 15));
            fp16_to_real = (sign) ? -result : result;
        end
    end
    endfunction
    
    // Function to check if two FP16 values are equal (including special cases)
    function bit fp16_equal;
        input [15:0] a;
        input [15:0] b;
    begin
        // For NaN, we consider any NaN equal to any other NaN
        if ((a[14:10] == 5'b11111 && a[9:0] != 10'b0) && 
            (b[14:10] == 5'b11111 && b[9:0] != 10'b0)) begin
            fp16_equal = 1'b1;
        end else begin
            fp16_equal = (a === b);
        end
    end
    endfunction
    
    // Test task
    task test_multiply;
        input [15:0] a_val;
        input [15:0] b_val;
        input [15:0] expected;
        string test_name;
        real a_real, b_real, expected_real, result_real;
    begin
        test_count = test_count + 1;
        
        // Wait for ready
        @(posedge clk);
        while (!ready) @(posedge clk);
        
        // Apply inputs
        A = a_val;
        B = b_val;
        valid = 1'b1;
        
        // Wait for computation
        @(posedge clk);
        valid = 1'b0;
        
        // Wait for result (3 cycles for pipeline)
        repeat(3) @(posedge clk);
        
        // Convert to real for debug
        a_real = fp16_to_real(a_val);
        b_real = fp16_to_real(b_val);
        expected_real = fp16_to_real(expected);
        result_real = fp16_to_real(C);
        
        // Check result using fp16_equal function for proper NaN handling
        if (fp16_equal(C, expected)) begin
            pass_count = pass_count + 1;
            $display("PASS: Test %0d - A=%h, B=%h, Got=%h, Expected=%h", 
                     test_count, a_val, b_val, C, expected);
        end else begin
            fail_count = fail_count + 1;
            $display("FAIL: Test %0d - A=%h, B=%h, Got=%h, Expected=%h", 
                     test_count, a_val, b_val, C, expected);
            $display("      A_real=%f, B_real=%f, Got_real=%f, Expected_real=%f",
                     a_real, b_real, result_real, expected_real);
        end
    end
    endtask
    
    initial begin
        // Initialize
        clk = 0;
        rst_n = 0;
        valid = 0;
        A = 0;
        B = 0;
        
        // Generate waveform
        $dumpfile("Out/wave.vcd");
        $dumpvars(0, fp_multiplier_tb);
        
        // Reset
        #10 rst_n = 1;
        
        $display("Starting FP16 Multiplier Tests...");
        
        // Test 1: Basic Multiplication - Positive numbers
        test_multiply(FP16_TWO, FP16_THREE, FP16_SIX); // 2.0 * 3.0 = 6.0
        
        // Test 2: Negative numbers  
        test_multiply(FP16_NEG_TWO, FP16_THREE, 16'hC600); // -2.0 * 3.0 = -6.0
        
        // Test 3: Mixed signs
        test_multiply(FP16_NEG_TWO, FP16_NEG_THREE, FP16_SIX); // -2.0 * -3.0 = 6.0
        
        // Test 4: Zero multiplication
        test_multiply(FP16_ZERO, FP16_TWO, FP16_ZERO); // 0.0 * 2.0 = 0.0
        test_multiply(FP16_TWO, FP16_ZERO, FP16_ZERO); // 2.0 * 0.0 = 0.0
        
        // Test 5: Infinity multiplication
        test_multiply(FP16_INF, FP16_TWO, FP16_INF); // inf * 2.0 = inf
        test_multiply(FP16_NEG_INF, FP16_THREE, FP16_NEG_INF); // -inf * 3.0 = -inf
        
        // Test 6: NaN multiplication
        test_multiply(FP16_NAN, FP16_TWO, FP16_NAN); // NaN * 2.0 = NaN
        
        // Test 7: Infinity * Zero = NaN
        test_multiply(FP16_INF, FP16_ZERO, FP16_NAN);
        test_multiply(FP16_ZERO, FP16_INF, FP16_NAN);
        
        // Test 8: Subnormal numbers - expect underflow to zero
        test_multiply(FP16_MIN, FP16_MIN, FP16_ZERO); // Very small * very small = 0 (underflow)
        
        // Test 9: Overflow
        test_multiply(FP16_MAX, FP16_TWO, FP16_INF); // 65504 * 2.0 = inf (overflow)
        
        // Test 10: Edge case - maximum representable
        test_multiply(FP16_MAX, FP16_ONE, FP16_MAX); // 65504 * 1.0 = 65504
        
        // Additional test: 1.0 * 1.0 = 1.0
        test_multiply(FP16_ONE, FP16_ONE, FP16_ONE);
        
        // Final summary
        $display("\nTest Summary:");
        $display("Total Tests: %0d", test_count);
        $display("Passed: %0d", pass_count);
        $display("Failed: %0d", fail_count);
        
        if (fail_count == 0) begin
            $display("TestSuccessed");
        end else begin
            $display("TestFailed");
        end
        
        $finish;
    end

endmodule
\end{lstlisting}

\newpage
\section{功能六：时序和面积优化}
除了以上的功能，我还尝试了一下Explore QoR功能，对刚刚生成的浮点乘法器做面积、时序优化。

\begin{figure}[H]
    \centering
    \includegraphics[width=\textwidth]{面积优化.png}
    \caption{时序和面积优化}
\end{figure}

最优时序和最优面积的代码如下：
\begin{lstlisting}[frame=shadowbox,caption=vRTL251119114743\_779.v]
module fp_multiplier (
    input clk,
    input rst_n,
    input [15:0] A,
    input [15:0] B,
    input Ready,
    output reg [15:0] C,
    output reg Valid
);

    // Internal registers and wires
    reg [15:0] computed_C;
    
    // Input decomposition
    wire sign_A;
    wire [4:0] exp_A;
    wire [9:0] mant_A;
    wire sign_B;
    wire [4:0] exp_B;
    wire [9:0] mant_B;
    
    // Special case detection
    wire is_nan_A;
    wire is_inf_A;
    wire is_zero_A;
    wire is_nan_B;
    wire is_inf_B;
    wire is_zero_B;
    
    // Computation signals - stage 1
    wire [10:0] sig_A;
    wire [10:0] sig_B;
    wire [21:0] product_sig;
    wire [5:0] exp_A_adj;
    wire [5:0] exp_B_adj;
    wire [6:0] exp_sum;
    wire [6:0] exp_prod;
    
    // Normalization signals - stage 2
    wire [4:0] shift_amount;
    wire [21:0] normalized_sig;
    wire [6:0] exp_after_norm;
    
    // Rounding signals - stage 3
    wire [9:0] mant_part;
    wire G;
    wire R;
    wire S;
    wire round_up;
    wire [10:0] rounded_mant;
    wire [9:0] final_mant;
    wire [6:0] final_exp;
    
    // Pipeline registers
    reg [1:0] special_case_stage1;
    reg sign_prod_stage1;
    reg [21:0] product_sig_stage1;
    reg [6:0] exp_prod_stage1;
    
    reg [1:0] special_case_stage2;
    reg sign_prod_stage2;
    reg [21:0] normalized_sig_stage2;
    reg [6:0] exp_after_norm_stage2;
    
    // Assign input decomposition
    assign sign_A = A[15];
    assign exp_A = A[14:10];
    assign mant_A = A[9:0];
    assign sign_B = B[15];
    assign exp_B = B[14:10];
    assign mant_B = B[9:0];
    
    // Special case detection
    assign is_nan_A = (exp_A == 5'b11111) && (mant_A != 10'b0);
    assign is_inf_A = (exp_A == 5'b11111) && (mant_A == 10'b0);
    assign is_zero_A = (exp_A == 5'b0) && (mant_A == 10'b0);
    assign is_nan_B = (exp_B == 5'b11111) && (mant_B != 10'b0);
    assign is_inf_B = (exp_B == 5'b11111) && (mant_B == 10'b0);
    assign is_zero_B = (exp_B == 5'b0) && (mant_B == 10'b0);
    
    // Stage 1: Basic computation
    assign sig_A = (exp_A == 5'b0) ? {1'b0, mant_A} : {1'b1, mant_A};
    assign sig_B = (exp_B == 5'b0) ? {1'b0, mant_B} : {1'b1, mant_B};
    assign product_sig = sig_A * sig_B;
    assign exp_A_adj = (exp_A == 5'b0) ? 6'b100001 : {1'b0, exp_A};
    assign exp_B_adj = (exp_B == 5'b0) ? 6'b100001 : {1'b0, exp_B};
    assign exp_sum = exp_A_adj + exp_B_adj;
    assign exp_prod = exp_sum - 7'b0011110;
    
    // Stage 1 pipeline register
    always @(posedge clk or negedge rst_n) begin
        if (!rst_n) begin
            special_case_stage1 <= 2'b00;
            sign_prod_stage1 <= 1'b0;
            product_sig_stage1 <= 22'b0;
            exp_prod_stage1 <= 7'b0;
        end else if (Ready) begin
            // Encode special cases
            if (is_nan_A || is_nan_B) begin
                special_case_stage1 <= 2'b01; // NaN
            end else if ((is_inf_A && is_zero_B) || (is_zero_A && is_inf_B)) begin
                special_case_stage1 <= 2'b01; // NaN
            end else if (is_inf_A || is_inf_B) begin
                special_case_stage1 <= 2'b10; // Infinity
            end else if (is_zero_A || is_zero_B) begin
                special_case_stage1 <= 2'b11; // Zero
            end else begin
                special_case_stage1 <= 2'b00; // Normal case
            end
            
            sign_prod_stage1 <= sign_A ^ sign_B;
            product_sig_stage1 <= product_sig;
            exp_prod_stage1 <= exp_prod;
        end
    end
    
    // Stage 2: Normalization
    assign shift_amount = (product_sig_stage1[21]) ? 5'b0 :
                         (product_sig_stage1[20]) ? 5'b1 :
                         (product_sig_stage1[19]) ? 5'b10 :
                         (product_sig_stage1[18]) ? 5'b11 : 5'b0;
                         
    assign normalized_sig = product_sig_stage1 << shift_amount;
    assign exp_after_norm = exp_prod_stage1 - {2'b0, shift_amount};
    
    // Stage 2 pipeline register
    always @(posedge clk or negedge rst_n) begin
        if (!rst_n) begin
            special_case_stage2 <= 2'b00;
            sign_prod_stage2 <= 1'b0;
            normalized_sig_stage2 <= 22'b0;
            exp_after_norm_stage2 <= 7'b0;
        end else begin
            special_case_stage2 <= special_case_stage1;
            sign_prod_stage2 <= sign_prod_stage1;
            normalized_sig_stage2 <= normalized_sig;
            exp_after_norm_stage2 <= exp_after_norm;
        end
    end
    
    // Stage 3: Rounding and final computation
    assign mant_part = normalized_sig_stage2[20:11];
    assign G = normalized_sig_stage2[10];
    assign R = normalized_sig_stage2[9];
    assign S = |normalized_sig_stage2[8:0];
    
    assign round_up = (G && (R || S)) ? 1'b1 :
                     (G && !(R || S) && mant_part[0]) ? 1'b1 : 1'b0;
                     
    assign rounded_mant = {1'b0, mant_part} + round_up;
    assign final_mant = rounded_mant[10] ? rounded_mant[10:1] : rounded_mant[9:0];
    assign final_exp = rounded_mant[10] ? (exp_after_norm_stage2 + 7'b1) : exp_after_norm_stage2;
    
    // Final result computation
    always @(*) begin
        computed_C = 16'b0;
        
        case (special_case_stage2)
            2'b01: computed_C = {1'b0, 5'b11111, 10'b1}; // NaN
            2'b10: computed_C = {sign_prod_stage2, 5'b11111, 10'b0}; // Infinity
            2'b11: computed_C = {sign_prod_stage2, 5'b0, 10'b0}; // Zero
            default: begin
                if (final_exp[6] || final_exp[5:0] > 6'b001111) begin
                    computed_C = {sign_prod_stage2, 5'b11111, 10'b0};
                end else if (final_exp[5:0] < 6'b000001) begin
                    computed_C = {sign_prod_stage2, 5'b0, 10'b0};
                end else begin
                    computed_C = {sign_prod_stage2, final_exp[4:0], final_mant};
                end
            end
        endcase
    end
    
    // Output register
    always @(posedge clk or negedge rst_n) begin
        if (!rst_n) begin
            C <= 16'b0;
            Valid <= 1'b0;
        end else begin
            C <= computed_C;
            Valid <= Ready; // Valid follows Ready with pipeline delay
        end
    end

endmodule
\end{lstlisting}

\begin{lstlisting}[frame=shadowbox,caption=fp\_multiplier\_generated3.sv]
module fp_multiplier (
    input clk,
    input rst_n,
    input [15:0] A,
    input [15:0] B,
    input Ready,
    output reg [15:0] C,
    output reg Valid
);

    // Internal registers and wires
    reg [15:0] computed_C;
    
    // Input decomposition
    wire sign_A = A[15];
    wire [4:0] exp_A = A[14:10];
    wire [9:0] mant_A = A[9:0];
    wire sign_B = B[15];
    wire [4:0] exp_B = B[14:10];
    wire [9:0] mant_B = B[9:0];
    
    // Special case detection
    wire is_nan_A = (exp_A == 5'b11111) && (mant_A != 10'b0);
    wire is_inf_A = (exp_A == 5'b11111) && (mant_A == 10'b0);
    wire is_zero_A = (exp_A == 5'b0) && (mant_A == 10'b0);
    wire is_nan_B = (exp_B == 5'b11111) && (mant_B != 10'b0);
    wire is_inf_B = (exp_B == 5'b11111) && (mant_B == 10'b0);
    wire is_zero_B = (exp_B == 5'b0) && (mant_B == 10'b0);
    
    // Shared computation signals
    wire [10:0] sig_A = (exp_A == 5'b0) ? {1'b0, mant_A} : {1'b1, mant_A};
    wire [10:0] sig_B = (exp_B == 5'b0) ? {1'b0, mant_B} : {1'b1, mant_B};
    wire [21:0] product_sig = sig_A * sig_B;
    wire [5:0] exp_A_adj = (exp_A == 5'b0) ? 6'b100001 : {1'b0, exp_A};
    wire [5:0] exp_B_adj = (exp_B == 5'b0) ? 6'b100001 : {1'b0, exp_B};
    wire [6:0] exp_sum = exp_A_adj + exp_B_adj;
    wire [6:0] exp_prod = exp_sum - 7'b0011110;
    
    // Area-optimized normalization using shared logic
    wire [4:0] shift_amount;
    wire [21:0] normalized_sig;
    
    // Optimized shift amount calculation using priority encoder style
    assign shift_amount = (~product_sig[21] & ~product_sig[20] & ~product_sig[19] & ~product_sig[18]) ? 5'b0 :
                         (~product_sig[21] & ~product_sig[20] & ~product_sig[19] & product_sig[18]) ? 5'b11 :
                         (~product_sig[21] & ~product_sig[20] & product_sig[19]) ? 5'b10 :
                         (~product_sig[21] & product_sig[20]) ? 5'b1 : 5'b0;
    
    // Shared barrel shifter
    assign normalized_sig = product_sig << shift_amount;
    
    wire [6:0] exp_after_norm = exp_prod - {2'b0, shift_amount};
    
    // Shared rounding logic
    wire [9:0] mant_part = normalized_sig[20:11];
    wire G = normalized_sig[10];
    wire R = normalized_sig[9];
    wire S = |normalized_sig[8:0];
    wire round_up = (G && (R || S)) || (G && !(R || S) && mant_part[0]);
    wire [10:0] rounded_mant = {1'b0, mant_part} + round_up;
    
    wire [9:0] final_mant = rounded_mant[10] ? rounded_mant[10:1] : rounded_mant[9:0];
    wire [6:0] final_exp = rounded_mant[10] ? (exp_after_norm + 7'b1) : exp_after_norm;
    
    // Combined result computation
    always @(*) begin
        computed_C = 16'b0;
        
        if (is_nan_A || is_nan_B) begin
            computed_C = {1'b0, 5'b11111, 10'b1};
        end else if ((is_inf_A && is_zero_B) || (is_zero_A && is_inf_B)) begin
            computed_C = {1'b0, 5'b11111, 10'b1};
        end else if (is_inf_A || is_inf_B) begin
            computed_C = {sign_A ^ sign_B, 5'b11111, 10'b0};
        end else if (is_zero_A || is_zero_B) begin
            computed_C = {sign_A ^ sign_B, 5'b0, 10'b0};
        end else begin
            if (final_exp[6] || final_exp[5:0] > 6'b001111) begin
                computed_C = {sign_A ^ sign_B, 5'b11111, 10'b0};
            end else if (final_exp[5:0] < 6'b000001) begin
                computed_C = {sign_A ^ sign_B, 5'b0, 10'b0};
            end else begin
                computed_C = {sign_A ^ sign_B, final_exp[4:0], final_mant};
            end
        end
    end
    
    // Output register
    always @(posedge clk or negedge rst_n) begin
        if (!rst_n) begin
            C <= 16'b0;
            Valid <= 1'b0;
        end else begin
            if (Ready) begin
                C <= computed_C;
                Valid <= 1'b1;
            end else begin
                Valid <= 1'b0;
            end
        end
    end

endmodule
\end{lstlisting}

从代码可以看到，对于优化时序，其采用了四级流水线设计；而对于优化面积，其降低了时序逻辑和组合逻辑的个数。
\section{对合见工软UDA的使用评价，修改建议}
做完整个任务后，我深深感受到了这个UDA工具的强大。其功能十分完善，纠错、补全代码、QoR探索、一站式验证调试等已经能满足大部分
设计者基本的设计任务，同时使用手册还提供了更加高阶的使用操作。并且内嵌AI模型，可以在设计中的Chat聊天框中和AI对话，非常的方便且灵活。

除此之外，一些细节的设计也非常到位。比如在某个代码中做了改错的操作，会在代码开头显示过对这代码有过UDA的什么操作；在Chat聊天框中可以一键插入代码、保存代码、新旧代码对比，非常好用。

至于修改建议，可以在设置中的prompt add里添加一个直接读取文件的功能；同时在修改代码时，可以自动跑一遍仿真，将仿真结果作为AI下一次修改意见；以及在FILES区域
添加文件的时候，可能是有小bug，有几次点进去添加文件，结果文件夹中的代码文件根本没有显示。

\end{document}